% SPDX-License-Identifier: MIT
\chapter{The Python API}
\label{ch:python}

The \cmd{makeyourtree} package imports no GUI toolkit, so everything in this
chapter runs in a script, a notebook, a container or a CI job.

\section{Loading and inspecting}

\begin{lstlisting}[language=Python]
from makeyourtree.io import load_tree, save_tree

tree = load_tree("examples/primates.nwk")

print(len(tree.nodes))
print([n.name for n in tree.nodes if not n.children])
print(tree.rooted)
\end{lstlisting}

\cmd{load\_tree} detects the format from the content. Pass \cmd{format=} to
override it. \cmd{load\_trees} reads a multi-tree file.

\subsection{Diagnostics}

Parsing problems are collected rather than raised. To see them, pass a sink:

\begin{lstlisting}[language=Python]
from makeyourtree.core.diagnostics import DiagnosticSink

sink = DiagnosticSink()
tree = load_tree("messy.nwk", sink=sink)
for d in sink:
    print(d.code, d.message)
\end{lstlisting}

\section{Traversal}

\begin{quote}
\textbf{Never recurse over a tree.} A deeply unbalanced tree of 100{,}000 tips
has a path length of 100{,}000, and Python's recursion limit is 1{,}000. Use the
iterative traversals in \cmd{makeyourtree.core.traversal}; every traversal inside
\app{} is iterative for exactly this reason.
\end{quote}

\begin{lstlisting}[language=Python]
from makeyourtree.core import traversal

for node in traversal.preorder(tree.root):
    ...
for node in traversal.postorder(tree.root):
    ...
\end{lstlisting}

\section{Operations}

Operations return a \emph{command} rather than mutating immediately, which is
what makes them undoable in the application. From a script, apply it at once:

\begin{lstlisting}[language=Python]
from makeyourtree.ops import (midpoint_root, outgroup_root, ladderize,
                              rotate, collapse, prune, extract_subtree, unroot)

midpoint_root(tree).apply(tree)
ladderize(tree, ascending=True).apply(tree)
\end{lstlisting}

To keep an undo history, use a \cmd{CommandStack}:

\begin{lstlisting}[language=Python]
from makeyourtree.ops import CommandStack

stack = CommandStack()
stack.do(midpoint_root(tree), tree)
stack.do(ladderize(tree), tree)
stack.undo(tree)          # back to midpoint-rooted only
\end{lstlisting}

\begin{quote}
\textbf{A branch length belongs to an edge, not to a node.} It is stored on the
child because that is where a parent-pointer tree can hold it in constant time,
but it describes the edge \emph{parent \textrightarrow{} child}. Any operation
that reverses an edge --- rerooting, unrooting, grafting --- must carry the
length and support across the reversal. \app{} does; if you write your own tree
surgery, this is the invariant to test, and testing that pairwise patristic
distances survive a reroot is the way to catch a breach.
\end{quote}

\section{Composing a figure}

\begin{lstlisting}[language=Python]
from makeyourtree.doc import Document
from makeyourtree.layout import LayoutParams, LayoutMode, BranchMode
from makeyourtree.scene import compose
from makeyourtree.render import render_svg
from makeyourtree.style import LIGHT, DARK

doc = Document(
    tree=tree,
    params=LayoutParams(mode=LayoutMode.CIRCULAR, arc=340.0,
                        align_tips=True, row_spacing=16.0),
    theme=LIGHT,
)

scene = compose(doc)
with open("figure.svg", "w", encoding="utf-8") as fh:
    fh.write(render_svg(scene))
\end{lstlisting}

\cmd{compose} returns a \cmd{Scene}: plain data describing layered shapes and
text, with no reference to any rendering technology. \cmd{scene.count()} gives
the number of marks, and \cmd{scene.iter\_marks()} walks them in painting order.

\section{Attaching tracks}

From a file:

\begin{lstlisting}[language=Python]
from makeyourtree.annot import load_annotation

track = load_annotation("examples/primates_region.mytrack", tree)
doc.tracks.append(track)
\end{lstlisting}

Or built directly, which is the point of the API --- annotation computed from
data you already have, without writing a file:

\begin{lstlisting}[language=Python]
from makeyourtree.tracks.strip import ColorStripTrack

track = ColorStripTrack(
    title="Region",
    options={"thickness": 18,
             "colors": {"Africa": "#117733", "Asia": "#332288"}},
)
track.bind(tree, {"Homo_sapiens": ["Africa"],
                  "Macaca_mulatta": ["Asia"]}, columns=["region"])
doc.tracks.append(track)
\end{lstlisting}

\cmd{bind} takes a mapping of node name to a list of values, one per column, and
resolves the names against the tree. Chapter~\ref{ch:tracks} lists the class and
options for each track type.

\section{Raster and PDF output}

The core writes SVG. PNG and PDF come from the studio package, which requires Qt
and a \cmd{QApplication}:

\begin{lstlisting}[language=Python]
from PySide6.QtWidgets import QApplication
from makeyourtree_studio.export.raster import export_png, export_pdf

app = QApplication.instance() or QApplication([])
export_pdf(scene, "figure.pdf")
export_png(scene, "figure.png", scale=3.0)
\end{lstlisting}

\begin{quote}
\textbf{Fonts and the offscreen platform.} Qt's \cmd{offscreen} platform reports
zero font families, so text rasterises as empty boxes. If you are rendering PNGs
in a headless job and the text is missing, that is the cause. Clear
\cmd{QT\_QPA\_PLATFORM} before importing Qt, or use a virtual framebuffer such as
\cmd{xvfb-run}. SVG output from the core is unaffected, because it embeds text as
text.
\end{quote}

\section{Saving projects}

\begin{lstlisting}[language=Python]
from makeyourtree.doc.io import save_document, load_document

save_document(doc, "figure.mytree")
doc2 = load_document("figure.mytree")
\end{lstlisting}

\begin{quote}
\textbf{Do not put live objects in \cmd{document.metadata}.} It is serialised on
save, and an unserialisable value there will fail the write. Runtime-only keys
are stripped on save; anything else you add must be JSON-representable.
\end{quote}

\section{A complete example}

\begin{lstlisting}[language=Python]
from makeyourtree.io import load_tree
from makeyourtree.doc import Document
from makeyourtree.layout import LayoutParams, LayoutMode
from makeyourtree.ops import midpoint_root, ladderize
from makeyourtree.scene import compose
from makeyourtree.render import render_svg
from makeyourtree.annot import load_annotation

tree = load_tree("examples/bacteria.nwk")
midpoint_root(tree).apply(tree)
ladderize(tree).apply(tree)

doc = Document(tree=tree,
               params=LayoutParams(mode=LayoutMode.CIRCULAR,
                                   arc=340.0, align_tips=True))
for path in ("examples/bacteria_resistance.mytrack",
             "examples/bacteria_expression.mytrack"):
    doc.tracks.append(load_annotation(path, tree))

with open("figure.svg", "w", encoding="utf-8") as fh:
    fh.write(render_svg(compose(doc)))
\end{lstlisting}

\cmd{tools/make\_gallery.py} in the repository is a longer worked example: it
builds every figure in this manual, exercising all thirteen track types and every
layout mode.
